\section{模型评价与推广}

\subsection{模型的优点}

\textbf{一、以可行集替代点估计，结论带有保证而非仅有点值。} 本文全程不引入概率分布假设，而是维护"与已执行观测相容的世界"这一集合，并在其上取最坏值。由此得到的不是"目标大约在何处"，而是"目标必然在此区域内"；清除判据 $\varrho(P)\le20$ 也因此成为充要条件，而非经验阈值。对本题这种只给出误差界、不给出分布的设定，这是唯一能与题设严格对应的建模方式。

\textbf{二、把几何结论与后续决策直接打通。} 问题一输出的是区域直径、最小包围圆半径与保证清除位置集合这三个量，其中后两者正是问题三、四反复调用的接口。问题二输出的候选区域与最坏后验半径，则定义了"一次观测值不值得做"。整个模型体系不是四个独立的小问，而是一条从几何到决策的连续链，各问之间通过明确的数据接口传递结论。

\textbf{三、识别并正确更换了失效机制。} 第四问的建模不是给第三问"加一个源类型标签"，而是先逐项检查第三问的机制哪些仍然成立、哪些失效。本文明确指出无信号的语义（距离太远 vs. 处于背面）与覆盖终止证书这两项根本失效，并分别给出替代方案。这种"先说明旧方法为何失效、再给出修正"的组织方式，使模型扩展的理由是可检验的，而不是修辞性的。

\textbf{四、给出了可检验的终止依据。} 本题源数未知，终止条件是正确性的关键。本文用两类依据收口：达到 16 个的题目上界，或每个未发现频道都有全场覆盖证书。覆盖证书采用整格距离判定而非中心采样，第四问的证书则通过凸组合与凸包论证做到与朝向无关。这使"任务完成"成为一个可以被独立复核的命题，而不是程序自行宣布的状态。

\textbf{五、计费统一、结果可复现。} 全部策略共用式 \eqref{eq:time-ledger} 这一时间账本，使少走路、少检测、少切换、少失败尝试在同一尺度上可比；每次运行保存配置、场景、源码快照与散列、逐动作记录，数字可由脚本从原始结果重新生成。

\subsection{模型的局限性}

\textbf{一、第三、四问尚未完成题目要求的三次正式测试。} 这是本文最需要说明的限制。目前第三问的全部数值来自本地合成场景与一次官方演练，第四问的全部数值来自本地合成场景，尚无官方演练记录。因此表 \ref{tab:q3-formal} 与表 \ref{tab:q4-formal} 留空，相应的日志文件亦未产出。本地合成场景的生成分布（位置分布、半径分布、误差场、定向源比例）均为本文设定，不是题目给定的事实，其均值不能作为官方成绩的预期。

\textbf{二、策略的局部性。} 第三、四问的路线规划是滚动时域启发式：每步只执行当前路线第一项，评分为局部估计。第三问的路线 2-opt 与扫描点的删除/移动/重排是交替改进，整体组合问题非凸，不构成全局最优证明。第四问的下一测点评分式 \eqref{eq:q4-score} 中的权重与固定项均为工程参数，未经系统标定。独立审计也显示，局部成本的改善并没有稳定减少整场移动距离。

\textbf{三、尚无"任意合法场景必然完成"的证明。} 第三问的定位补测设有次数上限（每目标最多 16 次），超出预算即记录未完成并退出；有限预算下并未证明所有允许的误差函数都必然完成。第四问的完整性逻辑依赖于题面静态物理条件与观测误差界成立，且未被预算截断。

\textbf{四、数值实现为浮点加保守余量，而非形式化证明。} 所有几何计算使用浮点与显式容差（角域半宽额外加 0.0001$^\circ$，排除圆半径取 999.8 米或 999.99 米，清除判据取 19.5 米等）。这些余量使结果偏保守，但不构成区间算术意义下的严格证书；对极端病态的输入仍须检查状态与残差。

\textbf{五、效率提升的归因不完整。} 第四问 40 场配对中路线与可见性选点同时变化，当前没有足够消融证据把 26.00\% 的收益归于单一模块；第三问中失败信息更新的效率贡献经 20 场消融检验为不显著。此外，第三问的 v6 在固定极端与棋盘误差场下优势消失，说明其收益依赖误差的空间结构，不能推广到未知的官方误差场。

\textbf{六、尚未达到内部研究目标，但该目标本身也需正确理解。} 团队设定的"200 秒/点"目标目前未达成。圆周 10 源场景的移动下界为 233.64 秒/点，这否定了在所有合法场景上获得 200 秒硬保证的可能性，但并未证明特定分布下的均值不可达。需要强调，"200 秒/点"是团队内部的研究目标而非题目要求，不应与题目指标混淆。

\subsection{模型的推广}

\textbf{一、推广到更多观测站与三维情形。} 半平面表示与顶点枚举方法对观测数量没有本质限制，观测很多时可将 $O(m^3)$ 的交点枚举替换为 $O(m\log m)$ 的半平面交算法与旋转卡壳，用于直径计算。若测向机与干扰源存在不可忽略的高度差，可将二维半平面推广为三维中的锥面约束，$P$ 相应地成为多面锥的交集，直径与最小包围球的论证结构（凸组合 + 范数凸性）仍然成立，只是支撑点的个数上界由 3 变为 4。

\textbf{二、推广到移动源或接收半径变化的源。} 本文的正负观测半平面约束式 \eqref{eq:pos-neg-halfplane} 依赖"同一源接收半径固定"。若接收半径随时间变化，该约束失效，但圆盘排除式 \eqref{eq:disk-exclusion} 仍可使用（只需把 1000 米下界替换为当时的实际下界）。若源缓慢移动，则需把静态可行集升级为随时间演化的集合，并在其上引入运动模型。

\textbf{三、推广到更一般的朝向未知问题。} 第四问的凸包排除定理的适用条件与"定向源为半平面型可见性"这一具体形式无关，只要可见性条件可以写成"某个线性泛函非负"，该论证即可成立。因此该定理可直接推广到任意个数的定向观测站，也可推广到辐射范围不是恰好 $180^\circ$ 但包含某个半平面的情形。这一推广的关键在于保留"凸区域 $\subseteq$ 无信号点的凸包"这一结构，而不在于朝向参数的具体维度。

\textbf{四、推广到带概率信息的场景。} 本文有意克制地不使用概率假设。若实际系统中能通过大量重复检测标定出误差的统计规律（注意题目指出同一地点误差固定，故标定必须跨地点进行），则可在现有可行集之外附加置信区域，并把最坏值判据替换为机会约束或鲁棒—随机混合判据。此时本文的几何骨架仍可直接复用，只需替换决策目标，而不必重建整个模型。

\textbf{五、推广到多机协同。} 若同时派出多只机器狗，任务点集、覆盖证书与路线评分都需要共享状态；式 \eqref{eq:q4-score} 中的收益项需要改为"对整个团队的信息增益"，覆盖证书中的无信号记录则可以直接跨机合并（因为其有效性只依赖站点位置与频道，不依赖执行者）。这将是本模型最自然的下一步扩展方向之一。
